% ============================================================================
% Fichier cible : ARIEL_LOGMINER_MEMOIRE_FINAL/chapters/chapitre2_etat_art.tex
% Passage à remplacer : de « \subsection{CSE-CIC-IDS2018} » (inclus) à
%   « \subsection{HDFS} » (exclu).
% Agent : Agent 3 — État de l'art et positionnement, chantier 2026-11-05.
% Sources : Sharafaldin et al. ; source officielle CSE-CIC-IDS2018 ; Liu et al.
%   (CNS 2022, DOI 10.1109/CNS56114.2022.9947235).
% ============================================================================

\subsection{CSE-CIC-IDS2018}

CSE-CIC-IDS2018 prolonge la méthode de génération de CICIDS2017 dans une
infrastructure plus vaste et publie des captures, des journaux et des flux décrits par
des caractéristiques extraites avec CICFlowMeter
\cite{sharafaldin2018cicids,cic2018csecicids}. Sa taille et la diversité annoncée des
scénarios expliquent son emploi fréquent comme benchmark. Elles n'en font toutefois pas
un jeu indépendant de CICIDS2017 : les deux corpus partagent des principes de génération,
un extracteur de flux et plusieurs catégories d'attaque.

Liu et al. ont réexaminé CICIDS2017 et CSE-CIC-IDS2018 sur l'ensemble de leur chaîne de
création. Ils documentent des erreurs dans l'orchestration des attaques, la génération
des caractéristiques, la documentation et l'étiquetage, puis montrent que ces erreurs
peuvent déstabiliser les résultats expérimentaux \cite{liu2022error}. Leur conclusion ne
consiste pas à proscrire le corpus, mais à imposer trois précautions : identifier
précisément les fichiers utilisés, contrôler leur intégrité, et éviter qu'une variable
liée au dispositif de collecte serve de substitut à la classe.

Dans ce mémoire, CSE-CIC-IDS2018 est donc traité comme un objet d'audit de la validité
externe plutôt que comme une preuve autonome de transfert. Une performance élevée sur
deux journées ne devient interprétable qu'après inversion du sens de transfert, ablation
des variables dominantes et comparaison avec des règles mono-variable. Cette lecture
critique rattache directement le protocole du chapitre expérimental aux limites déjà
établies dans la littérature.

